% 本文件基于 mra.py 的实际实现整理，可作为毕业论文中的章节文件直接引用。

\chapter{模型介绍}
\label{chap:model_intro}

\section{模型设计背景与总体思路}

工业过程异常检测的核心任务是根据传感器采集到的多变量时间序列，及时识别设备故障、工况漂移或过程失稳等异常状态。与一般的分类任务不同，工业现场往往存在两个突出困难：一是不同传感器之间存在复杂耦合关系，单变量建模难以完整刻画系统动态；二是实际数据中常伴随缺失值、噪声以及多时间尺度变化，仅依靠单一时域特征容易造成异常模式刻画不充分。

结合上述特点，\texttt{mra.py} 中实现了一种自适应图--频率融合异常检测网络 AGF-ADNet（Adaptive Graph-Frequency Anomaly Detection Network）。该模型采用基于重构的异常检测思想，即先利用正常工况数据训练模型学习序列的内在演化规律，再通过测试阶段的重构误差衡量样本是否偏离正常模式。如果输入样本属于正常状态，则模型通常能够较准确地恢复其主要动态结构；若输入样本含有故障扰动或显著分布偏移，模型的重构误差将明显增大，从而为异常判别提供依据。

从实现机制上看，AGF-ADNet 不是一个单一路径的序列模型，而是一个融合了图结构学习、时域建模、频域增强、缺失值插补以及 Transformer 全局重构的复合网络。其设计逻辑可以概括为以下几点：
\begin{itemize}
    \item 使用自适应图学习模块刻画变量之间的动态关联关系，使模型能够识别“哪些传感器之间应当互相参考”；
    \item 使用图卷积与多尺度时序卷积联合提取时域信息，以同时建模变量间依赖和不同时间尺度的动态模式；
    \item 使用频域分支对序列的频谱结构进行增强，以补充时域建模对周期性、振荡性模式刻画不足的问题；
    \item 使用门控融合模块自适应整合时域与频域特征，并在缺失位置生成插补结果；
    \item 使用 Transformer 编码器在完整窗口范围内进行全局重构，最终输出每个时间步、每个变量的重构值。
\end{itemize}

因此，AGF-ADNet 的本质可以理解为一种“先补全、再重构、再以误差判别异常”的混合式深度模型。对于本科毕业设计而言，这样的结构既保留了较强的工程可实现性，又具有比较清晰的模块划分，便于从数学原理和代码实现两个角度展开论述。

\section{问题定义与符号说明}

\subsection{输入数据形式}

设原始工业过程数据表示为
\begin{equation}
    \mathbf{D} = [\mathbf{d}_1, \mathbf{d}_2, \ldots, \mathbf{d}_T]^\top \in \mathbb{R}^{T \times N},
\end{equation}
其中，$T$ 表示时间长度，$N$ 表示监测变量个数，$\mathbf{d}_t \in \mathbb{R}^{N}$ 表示第 $t$ 个时刻全部传感器组成的观测向量。

\texttt{mra.py} 中的数据通过读取多个 CSV 文件获得，每个文件不包含表头，列数即为变量个数。代码在运行时根据训练数据的列数自动确定模型中的节点数，因此虽然类定义中给出了默认值 $N=18$，但真正使用的变量维数由数据集本身决定。

\subsection{缺失掩码定义}

代码中采用如下掩码约定：
\begin{equation}
    M_{t,n} =
    \begin{cases}
        1, & \text{若第 } t \text{ 个时刻第 } n \text{ 个变量缺失}, \\
        0, & \text{若该位置有观测值}.
    \end{cases}
\end{equation}

这一定义与部分文献中“$1$ 表示已观测、$0$ 表示缺失”的写法正好相反，因此在撰写论文时必须特别说明，以免后续损失函数和插补公式的语义发生混淆。相应地，可定义观测掩码为
\begin{equation}
    O_{t,n} = 1 - M_{t,n},
\end{equation}
其中 $O_{t,n}=1$ 表示当前位置可用于监督或评分。

\subsection{窗口化输入张量}

为了便于神经网络建模，脚本采用长度为 $S$ 的滑动窗口构造样本。窗口长度在当前实现中固定为
\begin{equation}
    S = 60.
\end{equation}

在形成批量样本后，模型输入张量记为
\begin{equation}
    \mathbf{X} \in \mathbb{R}^{B \times S \times N}, \qquad
    \mathbf{M} \in \{0,1\}^{B \times S \times N},
\end{equation}
其中 $B$ 为批大小，默认取 $32$。张量元素 $X_{b,s,n}$ 表示第 $b$ 个样本在窗口内第 $s$ 个时间步、第 $n$ 个变量上的数值。

\subsection{模型输出}

对于一个批次的输入，AGF-ADNet 的前向传播可以表示为
\begin{equation}
    (\hat{\mathbf{X}}, \mathbf{A}, \mathbf{X}^{\mathrm{imp}})
    = f_{\theta}(\mathbf{X}^{\mathrm{in}}, \mathbf{M}),
\end{equation}
其中：
\begin{itemize}
    \item $\mathbf{X}^{\mathrm{in}}$ 为经过缺失位置清零后的实际输入；
    \item $\hat{\mathbf{X}} \in \mathbb{R}^{B \times S \times N}$ 为最终重构结果；
    \item $\mathbf{A} \in \mathbb{R}^{B \times N \times N}$ 为自适应学习到的邻接矩阵；
    \item $\mathbf{X}^{\mathrm{imp}} \in \mathbb{R}^{B \times S \times N}$ 为时频融合后得到的插补结果。
\end{itemize}

\section{数据预处理与样本构造}

\subsection{数据读取与拼接}

\texttt{DatasetBuilder.load\_dir()} 会读取指定目录下所有匹配模式的 CSV 文件，并按照文件名排序后在时间维进行拼接。设第 $k$ 个文件对应的矩阵为 $\mathbf{D}^{(k)} \in \mathbb{R}^{T_k \times N}$，则整体数据可写为
\begin{equation}
    \mathbf{D} =
    \begin{bmatrix}
        \mathbf{D}^{(1)} \\
        \mathbf{D}^{(2)} \\
        \vdots \\
        \mathbf{D}^{(K)}
    \end{bmatrix}.
\end{equation}

缺失掩码则由 \texttt{NaN} 位置直接生成：
\begin{equation}
    M_{t,n} = \mathbb{I}(D_{t,n}\ \text{is NaN}),
\end{equation}
其中 $\mathbb{I}(\cdot)$ 为指示函数。

\subsection{标准化处理}

标准化模块使用 \texttt{StandardScaler} 完成 Z-score 变换。由于标准化算子无法直接处理 \texttt{NaN}，代码先将缺失位置临时填充为 $0$，然后仅基于训练集估计均值与标准差。设经零填充后的训练数据为 $\bar{\mathbf{D}}^{\mathrm{tr}}$，则第 $n$ 个变量的标准化参数为
\begin{equation}
    \mu_n = \frac{1}{T_{\mathrm{tr}}}\sum_{t=1}^{T_{\mathrm{tr}}} \bar{D}^{\mathrm{tr}}_{t,n},
    \qquad
    \sigma_n = \sqrt{\frac{1}{T_{\mathrm{tr}}}\sum_{t=1}^{T_{\mathrm{tr}}}
    \left(\bar{D}^{\mathrm{tr}}_{t,n} - \mu_n\right)^2 }.
\end{equation}

随后训练集和测试集都使用同一组统计量进行缩放：
\begin{equation}
    D^{\mathrm{scaled}}_{t,n} = \frac{\bar{D}_{t,n} - \mu_n}{\sigma_n}.
\end{equation}

这种做法符合机器学习实验中的基本原则，即测试集不能参与标准化参数的估计，从而避免信息泄露。

\subsection{滑动窗口构造}

给定序列长度 $T$、窗口长度 $S$ 和步长 $\delta=1$，脚本以每一个时间步作为窗口终点构造样本。如果当前时刻历史长度不足 $S$，则使用首个样本进行前端填充。第 $i$ 个窗口可写为
\begin{equation}
    \mathbf{X}^{(i)} =
    \begin{cases}
        [\underbrace{\mathbf{d}_1, \ldots, \mathbf{d}_1}_{S-i-1}, \mathbf{d}_1, \mathbf{d}_2, \ldots, \mathbf{d}_{i+1}], & i < S, \\
        [\mathbf{d}_{i-S+2}, \ldots, \mathbf{d}_{i+1}], & i \geq S.
    \end{cases}
\end{equation}

对应的掩码窗口 $\mathbf{M}^{(i)}$ 按相同方式生成。这样处理有两个优点：其一，不会因为序列起始位置历史长度不足而丢弃早期样本；其二，可保证窗口构造后样本总数与原始时间长度一致，更便于后续分数序列与时间轴对齐。

\subsection{输入清零操作}

在模型前向传播前，脚本通过
\begin{equation}
    \mathbf{X}^{\mathrm{in}} = \mathbf{X} \odot (1 - \mathbf{M})
\end{equation}
将缺失位置直接置零，其中 $\odot$ 表示逐元素乘法。这样处理后，模型看到的是“观测值 + 缺失位置零占位”的输入张量，而缺失信息本身则通过掩码张量显式传入网络。

\section{模型总体结构}

从模块级别看，AGF-ADNet 的数据流可以概括为
\begin{align}
    \mathbf{A} &= \operatorname{GraphLearner}(\mathbf{X}^{\mathrm{in}}, \mathbf{M}), \\
    \mathbf{H}^{\mathrm{time}} &= \operatorname{TimeBranch}(\mathbf{X}^{\mathrm{in}}, \mathbf{A}), \\
    \mathbf{H}^{\mathrm{freq}} &= \operatorname{FreqBranch}(\mathbf{X}^{\mathrm{in}}), \\
    \mathbf{X}^{\mathrm{imp}} &= \operatorname{Fusion}(\mathbf{H}^{\mathrm{time}}, \mathbf{H}^{\mathrm{freq}}), \\
    \mathbf{X}^{\mathrm{filled}} &= \mathbf{X}^{\mathrm{in}} \odot (1-\mathbf{M}) + \mathbf{X}^{\mathrm{imp}} \odot \mathbf{M}, \\
    \hat{\mathbf{X}} &= \operatorname{TransformerRecon}(\mathbf{X}^{\mathrm{filled}}).
\end{align}

由此可见，模型并不是直接对原始输入进行重构，而是先通过图结构和时频信息融合得到一个插补结果，用它补足缺失位置，再对完整窗口进行全局重构。这样的两阶段设计能够显著增强模型对缺失场景和复杂动态场景的适应能力。

\section{自适应图学习模块}

\subsection{设计动机}

工业过程中多个变量之间往往存在显著耦合关系。例如，流量变化会影响压力，温度变化又可能进一步传导至反应速率或产品浓度。如果完全忽略这种变量间结构，仅依靠逐变量时序建模，模型会损失大量有效信息。因此，AGF-ADNet 首先使用图结构来表达变量之间的依赖关系，并且不依赖人工先验图，而是通过可学习方式自动生成邻接矩阵。

\subsection{动态上下文特征构造}

在 \texttt{GraphLearner} 中，每个变量节点都需要一个能够描述当前窗口状态的特征向量。代码从输入窗口中为每个节点提取四类动态统计量：
\begin{itemize}
    \item 观测均值；
    \item 观测标准差；
    \item 最近一次观测值；
    \item 缺失比例。
\end{itemize}

设第 $b$ 个样本中第 $n$ 个节点的窗口序列为 $\{x_{b,1,n}, \ldots, x_{b,S,n}\}$，观测掩码为 $o_{b,s,n}=1-m_{b,s,n}$，则上述统计量分别为
\begin{align}
    \mu_{b,n} &= \frac{\sum_{s=1}^{S} x_{b,s,n} o_{b,s,n}}
    {\sum_{s=1}^{S} o_{b,s,n} + \varepsilon}, \\
    \sigma_{b,n} &= \sqrt{
        \frac{\sum_{s=1}^{S}\left(x_{b,s,n} - \mu_{b,n}\right)^2 o_{b,s,n}}
        {\sum_{s=1}^{S} o_{b,s,n} + \varepsilon}
        + 10^{-6}}, \\
    \ell_{b,n} &= x_{b,s^\ast,n}, \qquad
    s^\ast = \arg\max_{s}(s \cdot o_{b,s,n}), \\
    r_{b,n} &= \frac{1}{S}\sum_{s=1}^{S} m_{b,s,n},
\end{align}
其中 $\varepsilon$ 为防止分母为零的极小常数。于是节点的动态描述向量可以写为
\begin{equation}
    \mathbf{d}_{b,n} = [\mu_{b,n}, \sigma_{b,n}, \ell_{b,n}, r_{b,n}] \in \mathbb{R}^{4}.
\end{equation}

上述设计具有明确的工程意义：均值和标准差用于刻画窗口整体分布，最近观测值强调当前状态，缺失比例则反映该节点在当前窗口中的信息可靠性。

\subsection{静态上下文与坐标先验}

除动态特征外，代码还为每个节点引入两类可学习静态参数：
\begin{itemize}
    \item \texttt{static\_context}：维度为 $d_s=8$ 的静态上下文向量；
    \item \texttt{node\_coords}：维度为 $d_c=8$ 的节点坐标向量。
\end{itemize}

其中，静态上下文向量可理解为节点的全局身份编码，用于补充长期稳定的结构信息；节点坐标则用于构造一个基于距离的先验邻接矩阵。若未提供外部基础邻接矩阵，则图先验定义为
\begin{equation}
    P_{i,j} =
    \begin{cases}
        \dfrac{1}{1+\|\mathbf{u}_i - \mathbf{u}_j\|_2}, & i \neq j, \\
        0, & i=j,
    \end{cases}
\end{equation}
其中 $\mathbf{u}_i \in \mathbb{R}^{d_c}$ 表示第 $i$ 个节点的可学习坐标。该公式的含义是：坐标越接近的节点，其潜在关联程度越高。

\subsection{节点表示编码}

将动态特征与静态上下文拼接后，可得到节点输入向量
\begin{equation}
    \mathbf{q}_{b,n} = [\mathbf{d}_{b,n}, \mathbf{c}_n] \in \mathbb{R}^{4+d_s},
\end{equation}
其中 $\mathbf{c}_n$ 为第 $n$ 个节点的静态上下文。随后，代码使用两层感知机进行编码：
\begin{equation}
    \mathbf{h}_{b,n} =
    \phi\left(
        \mathbf{W}_2
        \phi(\mathbf{W}_1 \mathbf{q}_{b,n} + \mathbf{b}_1)
        + \mathbf{b}_2
    \right),
\end{equation}
其中 $\phi(\cdot)$ 为 ReLU 激活函数，隐藏维度取 $32$。

\subsection{边权生成机制}

为了估计节点 $i$ 与节点 $j$ 之间的关联强度，代码为每一对节点构造如下配对特征：
\begin{equation}
    \mathbf{p}_{b,i,j} =
    [\mathbf{h}_{b,i},
    \mathbf{h}_{b,j},
    |\mathbf{h}_{b,i} - \mathbf{h}_{b,j}|,
    \mathbf{h}_{b,i} \odot \mathbf{h}_{b,j}],
\end{equation}
其中 $|\cdot|$ 表示逐元素绝对值，$\odot$ 表示逐元素乘法。该设计同时编码了节点自身表示、差异关系和交互关系，能够较充分地反映两节点之间的相似性与互补性。

边修正项由一个两层 MLP 给出：
\begin{equation}
    E_{b,i,j} = \operatorname{ReLU}\!\left(
        \mathbf{w}_e^\top
        \phi(\mathbf{W}_e \mathbf{p}_{b,i,j} + \mathbf{b}_e)
        + b_e'
    \right).
\end{equation}

最终的自适应邻接矩阵由“边修正项”和“图先验”逐元素相乘得到：
\begin{equation}
    A'_{b,i,j} = E_{b,i,j} \cdot P_{i,j}.
\end{equation}

随后，模型显式加入自环：
\begin{equation}
    A''_{b,i,j} = A'_{b,i,j} + \lambda \delta_{i,j},
\end{equation}
其中 $\lambda=1.0$，$\delta_{i,j}$ 为 Kronecker delta 函数。自环的作用是保留节点自身信息，避免图卷积过程中节点特征完全由邻居决定。

\subsection{对称归一化}

为稳定图卷积运算，代码对邻接矩阵进行对称归一化。定义度矩阵
\begin{equation}
    D_{b,i,i} = \sum_{j=1}^{N} A''_{b,i,j},
\end{equation}
则归一化邻接矩阵为
\begin{equation}
    \tilde{\mathbf{A}}_b = \mathbf{D}_b^{-1/2}\mathbf{A}''_b\mathbf{D}_b^{-1/2}.
\end{equation}

值得注意的是，该邻接矩阵是按批次动态生成的，即不同输入窗口可以对应不同的图结构。这一特性使模型具有更强的状态感知能力，也意味着图学习结果不是全局固定图，而是“样本相关图”。

\section{时域建模分支}

\subsection{图卷积层}

时域分支的第一部分是图卷积层 \texttt{GCNLayer}。其输入为
\begin{equation}
    \mathbf{X}^{\mathrm{in}} \in \mathbb{R}^{B \times S \times N}.
\end{equation}
代码先在最后增加一个特征维，得到
\begin{equation}
    \mathbf{X}^{4\mathrm{D}} \in \mathbb{R}^{B \times S \times N \times 1},
\end{equation}
然后对最后一维进行线性映射，再结合邻接矩阵做图传播：
\begin{equation}
    \mathbf{H}^{\mathrm{gcn}}_{b,s,:,:}
    = \tilde{\mathbf{A}}_b
    \left(
        \mathbf{X}^{4\mathrm{D}}_{b,s,:,:}\mathbf{W}_g + \mathbf{b}_g
    \right).
\end{equation}

由于输入维度和输出维度都设置为 $1$，因此这一层更接近“在变量关系图上做一次结构感知平滑与重分配”，其核心价值在于把某个传感器的表示与相关传感器的信息进行融合。

\subsection{多尺度时间卷积网络}

时域分支的第二部分是 \texttt{MultiScaleTCN}。该模块在每个变量通道上分别施加多种不同感受野的一维卷积，卷积核集合为
\begin{equation}
    \mathcal{K} = \{3,5,9\}.
\end{equation}

代码中采用 \texttt{groups=num\_nodes} 的卷积方式，因此每个变量独立进行时间卷积，属于典型的深度卷积（depthwise convolution）结构。设输入重新排列为
\begin{equation}
    \mathbf{X}^{\top} \in \mathbb{R}^{B \times N \times S},
\end{equation}
则对于第 $k \in \mathcal{K}$ 个尺度、第 $n$ 个变量，其卷积结果为
\begin{equation}
    y^{(k)}_{b,n,s}
    =
    \sum_{r=0}^{k-1}
    w^{(k)}_{n,r}
    x_{b,n,s+r-\lfloor k/2 \rfloor}
    + b^{(k)}_n.
\end{equation}

需要指出的是，\texttt{MultiScaleTCN} 类本身支持因果卷积和非因果卷积两种模式，但在 AGF-ADNet 主模型中设置为 \texttt{causal=False}，即采用对称填充。这样做意味着在窗口内部的重构阶段，模型可以同时利用过去和未来的上下文信息。由于本任务是窗口级重构而非在线一步预测，因此这种非因果设计是合理的。

不同尺度卷积输出会先沿新维度堆叠，再通过一个线性层进行自适应融合。若将三种尺度的结果记为 $\mathbf{Y}^{(3)}$、$\mathbf{Y}^{(5)}$ 和 $\mathbf{Y}^{(9)}$，则融合结果可表示为
\begin{equation}
    \mathbf{H}^{\mathrm{tcn}} =
    \mathbf{W}_f
    [\mathbf{Y}^{(3)}, \mathbf{Y}^{(5)}, \mathbf{Y}^{(9)}]
    + \mathbf{b}_f.
\end{equation}

\subsection{时域分支输出}

图卷积输出与多尺度卷积输出在代码中直接相加：
\begin{equation}
    \mathbf{H}^{\mathrm{time}} =
    \operatorname{LayerNorm}(\mathbf{H}^{\mathrm{gcn}})
    + \mathbf{H}^{\mathrm{tcn}}.
\end{equation}

该结构体现出一种“空间依赖 + 时间演化”联合建模思想：前者强调不同变量之间的信息传播，后者强调单变量在不同时间尺度上的局部动态，两者互补后形成时域分支的综合表示。

\section{频域建模分支}

\subsection{频域建模的必要性}

仅在时域中分析信号，虽然能够捕捉局部变化趋势，但对周期性波动、谐振模式和频率偏移等现象的刻画相对有限。工业过程中很多异常并不一定表现为绝对数值的剧烈变化，而可能表现为振荡频率变化、周期结构破坏或高频噪声增强。因此，在时域分支之外引入频域分支，有助于从另一个正交视角增强异常模式识别能力。

\subsection{快速傅里叶变换}

频域分支首先将输入从 $(B,S,N)$ 转换为 $(B,N,S)$，然后在时间维上执行实数快速傅里叶变换 \texttt{rFFT}：
\begin{equation}
    \mathbf{X}_f = \operatorname{rFFT}(\mathbf{X}^{\top})
    \in \mathbb{C}^{B \times N \times F},
\end{equation}
其中
\begin{equation}
    F = \left\lfloor \frac{S}{2} \right\rfloor + 1.
\end{equation}

由于输入是实值时间序列，采用 \texttt{rFFT} 可以避免对共轭对称频谱的重复计算，从而提高运算效率。

\subsection{频率特征表示}

设复数频谱可写为
\begin{equation}
    \mathbf{X}_f = \mathbf{R} + j\mathbf{I},
\end{equation}
其中 $\mathbf{R}$ 和 $\mathbf{I}$ 分别表示实部和虚部。代码虽然额外计算了频谱的幅值与相位，但真正送入神经网络的是实部和虚部的拼接表示：
\begin{equation}
    \mathbf{Z} = [\mathbf{R}; \mathbf{I}] \in \mathbb{R}^{B \times N \times 2F}.
\end{equation}

这种表示方式保留了完整的频谱信息，因为实部和虚部共同决定了幅值与相位。

\subsection{频率注意力机制}

为了突出更重要的频率分量，频域分支设计了一个两层感知机作为注意力网络：
\begin{equation}
    \boldsymbol{\alpha}
    = \sigma\!\left(
        \mathbf{W}^{\mathrm{att}}_2
        \phi(\mathbf{W}^{\mathrm{att}}_1\mathbf{Z} + \mathbf{b}^{\mathrm{att}}_1)
        + \mathbf{b}^{\mathrm{att}}_2
    \right)
    \in [0,1]^{B \times N \times F},
\end{equation}
其中 $\phi(\cdot)$ 为 ReLU 函数，$\sigma(\cdot)$ 为 Sigmoid 函数。由此得到的 $\alpha_{b,n,f}$ 表示第 $b$ 个样本中第 $n$ 个变量在第 $f$ 个频率分量上的重要程度。

\subsection{频谱增强与残差更新}

除注意力网络外，代码还构造了一个独立的频谱增强网络：
\begin{equation}
    \mathbf{Z}^{\mathrm{enh}}
    =
    \mathbf{W}^{\mathrm{enh}}_2
    \phi(\mathbf{W}^{\mathrm{enh}}_1\mathbf{Z} + \mathbf{b}^{\mathrm{enh}}_1)
    + \mathbf{b}^{\mathrm{enh}}_2
    \in \mathbb{R}^{B \times N \times 2F}.
\end{equation}

将其拆分为增强实部和增强虚部：
\begin{equation}
    \mathbf{Z}^{\mathrm{enh}} = [\Delta \mathbf{R}; \Delta \mathbf{I}],
\end{equation}
随后利用注意力权重进行逐频率调制：
\begin{equation}
    \tilde{\mathbf{R}} = \mathbf{R} + \boldsymbol{\alpha} \odot \Delta \mathbf{R},
    \qquad
    \tilde{\mathbf{I}} = \mathbf{I} + \boldsymbol{\alpha} \odot \Delta \mathbf{I}.
\end{equation}

因此增强后的复频谱为
\begin{equation}
    \tilde{\mathbf{X}}_f = \tilde{\mathbf{R}} + j\tilde{\mathbf{I}}.
\end{equation}

这种残差式更新而不是完全替换原频谱的做法，能够增强训练稳定性，并避免频域支路对原始谱结构造成过于激进的破坏。

\subsection{逆变换恢复时域表示}

最后，频域分支使用逆实数傅里叶变换将增强后的频谱恢复到时域：
\begin{equation}
    \mathbf{H}^{\mathrm{freq}}
    =
    \operatorname{irFFT}(\tilde{\mathbf{X}}_f)
    \in \mathbb{R}^{B \times N \times S}.
\end{equation}
再将其转回 $(B,S,N)$ 形式，并经过层归一化得到频域分支输出。由此，模型在同一输入窗口上同时获得时域表示与频域表示，为后续融合提供更全面的特征基础。

\section{时频门控融合与缺失值插补}

\subsection{门控融合机制}

时域分支和频域分支关注的信息不同，前者更擅长提取局部趋势与变量耦合，后者更擅长提取周期结构与频谱特征。因此，直接简单相加并不一定能够取得最佳效果。为此，代码采用门控融合模块 \texttt{GatedFusion} 自适应决定两者的组合比例。

设时域输出和频域输出分别为
\begin{equation}
    \mathbf{H}^{\mathrm{time}},\ \mathbf{H}^{\mathrm{freq}} \in \mathbb{R}^{B \times S \times N}.
\end{equation}
将它们在变量通道维上拼接后，得到
\begin{equation}
    \mathbf{C} \in \mathbb{R}^{B \times 2N \times S}.
\end{equation}

随后门控网络通过两层一维卷积产生门控系数：
\begin{equation}
    \mathbf{G}
    =
    \sigma\!\left(
        \operatorname{Conv}_2(
            \phi(\operatorname{Conv}_1(\mathbf{C}))
        )
    \right)
    \in [0,1]^{B \times N \times S}.
\end{equation}

将其转置回 $(B,S,N)$ 后，融合结果为
\begin{equation}
    \mathbf{X}^{\mathrm{imp}}
    =
    \mathbf{G} \odot \mathbf{H}^{\mathrm{time}}
    + (1-\mathbf{G}) \odot \mathbf{H}^{\mathrm{freq}}.
\end{equation}

上述公式说明：当某个位置的门控值接近 $1$ 时，模型更依赖时域特征；当门控值接近 $0$ 时，更依赖频域特征。门控融合后再经过 \texttt{LayerNorm}，以稳定特征尺度。

\subsection{缺失值插补}

门控融合结果随后用于缺失值插补。代码的插补公式为
\begin{equation}
    \mathbf{X}^{\mathrm{filled}}
    =
    \mathbf{X}^{\mathrm{in}} \odot (1-\mathbf{M})
    + \mathbf{X}^{\mathrm{imp}} \odot \mathbf{M}.
\end{equation}

该公式具有清晰的物理含义：对于已观测位置，直接保留原始输入；对于缺失位置，则使用模型给出的插补值。这样得到的 $\mathbf{X}^{\mathrm{filled}}$ 可以看作一个“观测值可信、缺失值可补”的完整窗口，为后续重构器提供更稳定的输入基础。

\section{Transformer 全局重构模块}

\subsection{输入投影与位置编码}

为了在时间维上进一步捕捉长程依赖关系，AGF-ADNet 在插补之后引入 Transformer 编码器。由于 Transformer 的输入特征维需要固定，因此代码先使用线性层将每个时间步的 $N$ 维变量向量投影到 $d=64$ 维隐空间：
\begin{equation}
    \mathbf{Z}_0 = \mathbf{X}^{\mathrm{filled}}\mathbf{W}_{\mathrm{in}} + \mathbf{b}_{\mathrm{in}},
    \qquad
    \mathbf{Z}_0 \in \mathbb{R}^{B \times S \times d}.
\end{equation}

在此基础上，再加上可学习位置编码 $\mathbf{P} \in \mathbb{R}^{1 \times S \times d}$：
\begin{equation}
    \mathbf{Z}_0' = \mathbf{Z}_0 + \mathbf{P}.
\end{equation}

与经典 Transformer 中固定的正弦位置编码不同，代码采用可训练位置参数，使模型能够在特定工业序列长度下自主学习更适合本任务的位置表示。

\subsection{Transformer 编码器}

脚本使用两层 \texttt{TransformerEncoderLayer} 构造编码器，每层包含多头自注意力机制和前馈网络，主要超参数如下：
\begin{itemize}
    \item 隐空间维度 $d=64$；
    \item 注意力头数 $h=4$；
    \item 前馈网络隐藏维度为 $128$；
    \item Dropout 为 $0.1$；
    \item 编码器层数为 $2$。
\end{itemize}

将第 $l$ 层编码器记为 $\mathcal{T}^{(l)}(\cdot)$，则其计算过程可以抽象表示为
\begin{equation}
    \mathbf{Z}_l = \mathcal{T}^{(l)}(\mathbf{Z}_{l-1}), \qquad l=1,2.
\end{equation}

由于代码未为 Transformer 显式设置因果掩码，因此它在窗口范围内可以访问完整时序上下文。这与前文的非因果 TCN 设计保持一致，说明模型的目标并不是在线单步预测，而是利用完整窗口进行高质量重构。

\subsection{输出投影与重构结果}

经过两层编码器后，脚本使用输出线性层将隐表示恢复到变量空间：
\begin{equation}
    \hat{\mathbf{X}} = \mathbf{Z}_2 \mathbf{W}_{\mathrm{out}} + \mathbf{b}_{\mathrm{out}},
    \qquad
    \hat{\mathbf{X}} \in \mathbb{R}^{B \times S \times N}.
\end{equation}

至此，模型完成从缺失输入到完整重构序列的映射。后续训练和异常检测都围绕 $\hat{\mathbf{X}}$ 与原始真实窗口 $\mathbf{X}$ 的差异展开。

\section{损失函数设计}

\subsection{重构损失}

AGF-ADNet 的训练目标由重构项、频域一致性项和图结构正则项三部分组成。首先，重构损失定义为
\begin{equation}
    \mathcal{L}_{\mathrm{rec}}
    =
    \frac{
        \left\|
            (\hat{\mathbf{X}} - \mathbf{X}) \odot \mathbf{T}
        \right\|_F^2
    }{
        \sum_{b,s,n} T_{b,s,n} + \varepsilon
    },
\end{equation}
其中 $\mathbf{T}$ 为目标监督掩码，表示哪些位置应参与当前批次的重构监督。

需要特别说明的是，$\mathbf{T}$ 并不总是简单等于观测掩码。在训练过程中，代码会在原本已观测的位置上随机再遮挡一部分值，以构造自监督恢复任务。因此，真正参与重构损失计算的通常是这些“人工遮挡位置”。

\subsection{频域一致性损失}

为了约束模型在频谱层面也保持重构一致性，脚本引入频域损失。定义仅在监督位置保留真实值、其余位置使用停止梯度后的重构值构成目标序列：
\begin{equation}
    \mathbf{X}^{\mathrm{freq\_target}}
    =
    \operatorname{stopgrad}(\hat{\mathbf{X}})\odot(1-\mathbf{T})
    + \mathbf{X}\odot\mathbf{T}.
\end{equation}

随后，在观测比例超过 $50\%$ 的样本上比较重构序列与目标序列的幅值频谱：
\begin{equation}
    \mathcal{L}_{\mathrm{freq}}
    =
    \frac{1}{|\Omega_v|}
    \sum_{b \in \Omega_v}
    \left\|
        \left|\operatorname{rFFT}(\hat{\mathbf{X}}_b)\right|
        -
        \left|\operatorname{rFFT}(\mathbf{X}^{\mathrm{freq\_target}}_b)\right|
    \right\|_F^2,
\end{equation}
其中 $\Omega_v$ 表示满足有效观测比例条件的样本集合。

该设计的意义在于：即使某些时域位置不直接参与监督，模型仍需要在整体频率结构上与目标信号保持一致，从而提高对周期和振荡模式的恢复能力。

\subsection{图结构熵正则项}

代码中所谓的 \texttt{sparsity\_loss} 实际上采用的是负熵形式：
\begin{equation}
    \mathcal{L}_{\mathrm{graph}}
    =
    -
    \frac{1}{BN}
    \sum_{b=1}^{B}
    \sum_{i=1}^{N}
    \sum_{j=1}^{N}
    \tilde{A}_{b,i,j}\log(\tilde{A}_{b,i,j}+\varepsilon).
\end{equation}

由于训练目标是最小化该项，因此优化过程会倾向于降低邻接分布的熵，使每个节点的连边更加集中，而不是平均地连接到所有节点。虽然它不是严格意义上的 $L_1$ 稀疏正则，但在效果上具有鼓励图结构“更尖锐、更有选择性”的作用。

\subsection{总损失函数}

综合上述三项，最终损失函数为
\begin{equation}
    \mathcal{L}
    =
    \mathcal{L}_{\mathrm{rec}}
    + \lambda_f \mathcal{L}_{\mathrm{freq}}
    + \lambda_g \mathcal{L}_{\mathrm{graph}},
\end{equation}
其中当前实现中
\begin{equation}
    \lambda_f = 0.1, \qquad \lambda_g = 0.1.
\end{equation}

该组合损失体现出多目标优化思想：既要求时域重构准确，也要求频域结构一致，同时还要求图结构具有较好的判别性和稳定性。

\section{训练策略与优化过程}

\subsection{自监督随机遮挡训练}

AGF-ADNet 采用自监督训练策略，其核心做法是在原本可观测的位置中随机挑选约 $10\%$ 进行再遮挡。设原始缺失掩码为 $\mathbf{M}$，随机遮挡得到的附加掩码为 $\mathbf{R}$，则训练输入掩码可表示为
\begin{equation}
    \mathbf{M}^{\mathrm{in}} = \mathbf{M} \lor \mathbf{R}.
\end{equation}

其中，$\lor$ 表示按位逻辑或。脚本随后根据 $\mathbf{M}^{\mathrm{in}}$ 构造输入
\begin{equation}
    \mathbf{X}^{\mathrm{in}} = \mathbf{X} \odot (1-\mathbf{M}^{\mathrm{in}}).
\end{equation}

与此同时，将随机遮挡位置作为主要监督目标，即
\begin{equation}
    \mathbf{T} = \mathbf{R}.
\end{equation}
如果某个批次中随机遮挡结果为空，则退化为使用全部已观测位置进行监督。这样设计的优点是无需额外人工标签，模型即可通过“遮挡--恢复”任务学习正常工况的表示。

\subsection{参数优化}

训练过程中，脚本使用 Adam 优化器，初始学习率为
\begin{equation}
    \eta_0 = 10^{-3}.
\end{equation}

同时使用 \texttt{ReduceLROnPlateau} 学习率调度器。当训练损失在若干轮内不再下降时，学习率按比例衰减：
\begin{equation}
    \eta \leftarrow 0.5 \eta.
\end{equation}

此外，为了抑制梯度爆炸，代码还使用最大范数为 $1.0$ 的梯度裁剪。当前脚本中的训练轮数为 $3$ 轮，这一数值更接近于示例实现或快速实验配置；若在正式毕业设计实验中追求更充分的收敛效果，通常还需要结合验证结果进一步调节训练轮数。

\subsection{最优模型保存}

脚本在训练过程中记录当前最小平均损失对应的模型参数：
\begin{equation}
    \theta^\ast = \arg\min_{\theta^{(e)}} \mathcal{L}^{(e)},
\end{equation}
其中 $\mathcal{L}^{(e)}$ 表示第 $e$ 轮训练的平均损失。训练结束后，程序会恢复到该最优参数状态，以降低偶然波动对测试阶段结果的影响。

\section{异常评分与判别机制}

\subsection{窗口级异常分数}

模型训练完成后，脚本会对每个窗口计算重构误差，并以此作为异常分数。设观测掩码为 $\mathbf{O}=1-\mathbf{M}$，则第 $b$ 个窗口的异常分数定义为
\begin{equation}
    s_b
    =
    \frac{
        \sum_{s=1}^{S}\sum_{n=1}^{N}
        \left(\hat{X}_{b,s,n} - X_{b,s,n}\right)^2 O_{b,s,n}
    }{
        \sum_{s=1}^{S}\sum_{n=1}^{N} O_{b,s,n} + \varepsilon
    }.
\end{equation}

从公式可以看出，异常评分只在真实观测位置上进行，缺失位置不会直接贡献测试误差。这一处理能够避免由于原始缺失值不可验证而造成评分失真。

\subsection{阈值设定}

脚本首先在训练集窗口上计算全部异常分数，然后用其均值加标准差作为阈值：
\begin{equation}
    \tau = \mu_{\mathrm{train}} + \sigma_{\mathrm{train}}.
\end{equation}

当测试窗口满足
\begin{equation}
    s_b > \tau
\end{equation}
时，即判定其为异常；否则判定为正常。该阈值方法实现简单、便于复现，适用于训练集主要由正常样本组成的情形。

\subsection{测试标签与评价指标}

需要指出的是，\texttt{mra.py} 中的测试标签并非从外部标签文件读取，而是通过 \texttt{build\_test\_labels()} 人为构造：默认将测试分数序列前半段视为正常，后半段视为异常。该设定显然依赖于当前数据组织方式，因此在论文中应将其说明为“当前实验脚本采用的标签构造规则”，而不宜笼统表述为通用异常检测标注机制。

在此基础上，代码计算了准确率（Accuracy）、精确率（Precision）、召回率（Recall）和 $F_1$ 值等分类指标，用于评估阈值判别效果。

\section{模型特点分析}

\subsection{模型优势}

结合代码实现，AGF-ADNet 主要具有以下几个特点：
\begin{itemize}
    \item \textbf{结构信息自适应学习。} 邻接矩阵无需人工指定，模型可根据当前窗口状态自动生成变量关系图；
    \item \textbf{时域与频域联合建模。} 模型同时利用局部时序模式和频谱结构，提高了对复杂异常模式的表达能力；
    \item \textbf{显式缺失值处理。} 通过掩码传递、门控融合和插补机制，模型能够在存在缺失值的情况下稳定工作；
    \item \textbf{全局上下文重构。} Transformer 编码器使模型能够在整个窗口范围内建模长程依赖，提升重构质量；
    \item \textbf{训练方式灵活。} 自监督随机遮挡机制减少了对人工标注的依赖，适合工业场景中“正常样本多、异常标签少”的实际情况。
\end{itemize}

\subsection{实现层面的注意事项}

从论文撰写角度看，还应注意以下几点：
\begin{itemize}
    \item 缺失掩码的定义为“$1$ 表示缺失”，这一点必须在符号表中明确；
    \item 图结构正则项实质上是低熵约束，而不是传统意义上的稀疏度约束；
    \item 主模型中的时序卷积和 Transformer 都是非因果的，因此模型更适合窗口重构场景，而不是严格实时预测场景；
    \item 当前脚本训练轮数较少，更多体现方法验证流程，正式实验时可进一步扩展训练与调参；
    \item 测试标签采用顺序划分构造，因此实验章节应把数据组织方式和评估前提写清楚。
\end{itemize}

\section{本章小结}

本章基于 \texttt{mra.py} 的实际实现，对 AGF-ADNet 的数据预处理流程、图结构学习机制、时域与频域特征提取方法、时频门控融合策略、Transformer 重构过程、损失函数设计以及异常评分方法进行了系统说明。整体来看，该模型围绕“自适应结构建模 + 时频联合表示 + 掩码重构检测”这一主线展开，既考虑了工业变量之间的耦合关系，也兼顾了缺失值场景和多尺度动态特征，是一种具有较强综合性的工业过程异常检测方法。

在后续实验章节中，可以进一步围绕该模型的检测性能、阈值选择、可视化结果以及与基线方法的对比情况展开分析，从而形成完整的本科毕业设计论文论证链条。
